\documentclass[12pt]{article}
\usepackage{amsmath}
\usepackage{amssymb}
\usepackage{graphicx}
\usepackage{hyperref}
\usepackage{geometry}
\geometry{a4paper, margin=1in}

\title{\textbf{Comparative Analysis of Regression Implementations: Rust (Burn) vs. PyTorch}}
\author{Technical Report}
\date{\today}

\begin{document}

\maketitle
\tableofcontents
\newpage

\section{Introduction}
This document outlines the architectural, mathematical, and deployment specifics of implementing a Neural Network-based Regression model across two disparate machine learning environments: Rust (utilizing the Burn framework) and Python (utilizing PyTorch). It covers the distinct model architecture decisions, dataset handling strategies, and specialized pipeline deployment techniques leveraging Network File Systems (NFS) mapping via Docker bounds.

\section{Model Architecture and Mathematical Formulation}

\subsection{Mathematical Foundation}
The core mathematical foundation deployed across both frameworks is a classical Feed-Forward Neural Network consisting of a single hidden dimension mapping inputs directly onto a continuous single-variable regression output.

For a given input feature vector $X \in \mathbb{R}^N$ (where $N$ dictates the feature size depending on the target dataset), the network's forward transformation can be represented sequentially as:
\begin{align}
    Z_1 &= X \cdot W_1^T + b_1 \quad &\text{(Input Projection)} \\
    A_1 &= \max(0, Z_1)        \quad &\text{(ReLU Activation)} \\
    \hat{Y} &= A_1 \cdot W_2^T + b_2 \quad &\text{(Output Projection)}
\end{align}

Where:
\begin{itemize}
    \item $W_1 \in \mathbb{R}^{H \times N}$ and $b_1 \in \mathbb{R}^H$ map the inputs onto the hidden vector space $H$.
    \item $\max(0, \cdot)$ denotes the Non-Linear Rectified Linear Unit (ReLU) mapping algorithm.
    \item $W_2 \in \mathbb{R}^{1 \times H}$ and $b_2 \in \mathbb{R}$ collapse the hidden abstraction onto the finalized regression scalar prediction $\hat{Y}$.
\end{itemize}

\subsection{Architectural Configurations}
While the mathematical foundations are identical, implementations slightly differ based on dataset selections within the modules:
\begin{itemize}
    \item \textbf{PyTorch Architecture:} Configures $N=13$ input features mapping to $H=64$ hidden parameters.
    \item \textbf{Rust (Burn) Architecture:} Configures $N=8$ input features concurrently mapping to $H=64$ hidden parameters.
\end{itemize}
In both configurations, standard parameter biases (`bias=True`) are included and automatically initialized.

\section{Training Pipelines}

Both codebases train the model iteratively tracking gradients via the Adam optimizer scaled against Mean Squared Error (MSE) loss logic:
\[ \text{MSE} = \frac{1}{B} \sum_{i=1}^{B} (Y_i - \hat{Y}_i)^2 \]

\subsection{PyTorch Context}
\begin{itemize}
    \item \textbf{Data Loading:} Automatically pulls the \textbf{Boston Housing} dataset array (.npz file) from an external Google API via `urllib` and manually partitions it down into an explicit $80/20$ split.
    \item \textbf{Telemetry Metrics:} Generates explicit hardware tracking loops inside the main epoch runner. Uses the `psutil` library to compute and stream epoch `iteration\_speed`, raw RAM consumption, and `cpu\_temp` hardware sensors parallel to the loss parameters.
\end{itemize}

\subsection{Rust (Burn) Context}
\begin{itemize}
    \item \textbf{Data Loading:} Links into Huggingface's dataset registry asynchronously targeting the \textbf{California Housing} SQLized splits mapping onto memory arrays via localized `HousingDistrictItem` structs.
    \item \textbf{Normalization Mapping:} Computes spatial min-max normalizations programmatically over inputs during training:
    \[ X_{norm} = \frac{X - \text{min}}{\text{max} - \text{min}} \]
    This logic restricts features within standard boundaries precluding exploding gradient derivations.
\end{itemize}

\section{Inference Pipeline and Docker NFS Integration}

Deploying these isolated pipelines necessitates radically different execution strategies, highlighting Python's heavyweight runtime dependency bottlenecks versus Rust's compile-time optimizations.

\subsection{PyTorch Inference Architecture}
Standard PyTorch Docker environments routinely eclipse several gigabytes due to CUDA bindings and generic scientific computation loops. To circumvent this inside microservices, the PyTorch inference pipeline mandates a hybrid Network File System (NFS) mapping architecture:
\begin{enumerate}
    \item \textbf{NFS Mounting (\texttt{mount\_libs.sh}):} Installs an external `nfs-common` client locally and binds the extensive python library volume from an external dedicated storage server (`172.16.203.14`) into the host machine's `/mnt/LSTM-libs` map.
    \item \textbf{Lightweight Container Image:} The backend \texttt{Dockerfile} avoids `pip install` commands completely, simply initializing a barebone `nvidia/cuda:12.1.1` image mapping Python $3.11$ system links. 
    \item \textbf{Volume Inject (\texttt{run\_container.sh}):} The script initializes the container enforcing `-v` flags that sync the NFS `/mnt/LSTM-libs` directory seamlessly onto the Docker's `/external-libs`. Crucially, it overrides the system \texttt{PYTHONPATH} to target those external `site-packages` at runtime.
    \item \textbf{Execution:} The `FastAPI` instance loads, bypasses massive disk pulls, links the models iteratively, and fields inbound `HousingFeatures` lists continuously.
\end{enumerate}

\subsection{Rust (Burn) Inference Architecture}
Rust handles Docker microservices inherently via statically linked deployments:
\begin{itemize}
    \item \textbf{Multi-Stage Compiling:} Executes a build phase operating within an oversized `rust:1.92-alpine` chain, ejecting the resulting binary onto an isolated stripped `alpine:3.23` environment structure.
    \item \textbf{Native Routing:} Utilizes \texttt{Axum} servers to establish the HTTP logic endpoints securely routing JSON payloads mapping to specific feature names (e.g. \texttt{median\_income}, \texttt{house\_age}).
\end{itemize}

\section{Language Specific Implementation Details}

\subsection{PyTorch-Specific Paradigms}
\begin{itemize}
    \item \textbf{Thread Clamping:} Due to inference optimization restrictions (especially running CPU variations alongside container structures), the `app.py` enforces explicit core binding calls via `torch.set\_num\_threads(1)` and `torch.set\_num\_interop\_threads(1)` securing computational resources and restricting OS context-switching overheads.
    \item \textbf{Matrix Array Verifications:} Manually inspects raw matrix vector mappings validating dimensions dynamically against numeric constraints: \texttt{len(x) != NUM\_FEATURES} triggering runtime panics before pipeline evaluations fail. 
    \item \textbf{Manual Hardware Moving:} The framework is heavily littered with required `.to(device)` mapping configurations switching inputs, datasets, targets, and models manually between the host and external components.
\end{itemize}

\subsection{Rust (Burn)-Specific Paradigms}
\begin{itemize}
    \item \textbf{Generic Compile-Time Shapes:} Dimension mappings and tensor validations are fundamentally enforced inside the Rust compiler boundaries via `<B, 2>` arrays indicating batches of distinct input structures mapping to `targets: Tensor<B, 1>`. Invalid sizes fail compilation, voiding the requirement for manual PyTorch matrix validations.
    \item \textbf{Struct Batching Protocols:} Inference doesn't evaluate primitive float arrays. Intead, the API relies on executing an overarching `HousingBatcher` which transforms specific struct domains (\texttt{HousingDistrictItem}) safely into tensor primitives while executing implicit `self.normalizer.to\_device(device)` logic silently against constants behind boundaries.
    \item \textbf{Record Deserialization:} States are strictly detached from models via standard `.mpk` maps. They invoke explicit \texttt{NoStdTrainingRecorder::new().load()} tracking traits unbinding memory limits inherent to standard dict serialization configurations natively loaded via `RegressionModelConfig`.
\end{itemize}

\end{document}
